\documentclass[11pt,a4paper]{article}

\usepackage[utf8]{inputenc}
\usepackage[T1]{fontenc}
\usepackage{geometry}
\usepackage{graphicx}
\usepackage{booktabs}
\usepackage{hyperref}
\usepackage{amsmath}
\usepackage{caption}
\usepackage{subcaption}
\usepackage{float}
\usepackage{xcolor}
\usepackage{listings}
\usepackage[numbers]{natbib}
\usepackage{enumitem}
\usepackage{multirow}

\geometry{margin=1in}

\lstset{
  basicstyle=\ttfamily\small,
  breaklines=true,
  frame=single,
  backgroundcolor=\color{gray!10},
  columns=fullflexible,
}

\definecolor{vtmaroon}{HTML}{861F41}
\definecolor{vtorange}{HTML}{E5751F}
\hypersetup{colorlinks=true,linkcolor=vtmaroon,citecolor=vtmaroon,urlcolor=vtorange}

\title{Porting CHIME to CXL-Attached Memory\\[0.5em]\large Final Report --- Part Two}
\author{Jason Cusati (\texttt{djjay@vt.edu})\\
CS 6204 --- Advanced Topics in Systems\\
Virginia Tech, Spring 2026}
\date{May 5, 2026}

\begin{document}

\maketitle

\begin{abstract}
%% TODO: 150-200 words. Cover:
%%  - CHIME's RDMA assumption
%%  - CXL as a point-to-point alternative
%%  - Our port: CxlTransport backend, USE_CXL flag, NUMA emulation
%%  - Headline finding (CXL vs RDMA latency/throughput on single-node)
%%  - Main contribution: open-source, reproducible CXL port of CHIME
CHIME (SOSP~'24) is a cache-efficient hybrid B+~tree/hopscotch index designed for RDMA-based disaggregated memory.
We describe a port of CHIME that replaces its RDMA transport with a CXL-attached memory abstraction, implemented via NUMA-node emulation on a dual-socket host.
We evaluate the CXL port against a single-node RDMA baseline on CloudLab r650 nodes using the YCSB workload suite and show that [RESULT~PLACEHOLDER].
Our port is minimally invasive: a single compile-time flag (\texttt{USE\_CXL=ON}) swaps the transport layer without changes to CHIME's index algorithm or concurrency protocol.
All code is available at \url{https://github.com/djjay0131/CHIME}.
\end{abstract}

%% ============================================================
\section{Introduction}
\label{sec:intro}
%% ============================================================

%% TODO: Motivation — why CXL for CHIME?
%%  - CXL 2.0/3.0 promises lower latency, higher bandwidth, cache coherence
%%  - Compare CXL (~200ns) to RDMA (~2us) load latency
%%  - CHIME's design assumes high per-op cost → interesting question whether its
%%    optimizations still pay off on CXL
%%  - This report: port, measure, and discuss

\subsection{Contributions}
\begin{itemize}[nosep]
    \item A minimally invasive port of CHIME to CXL-attached memory, using a compile-time transport switch.
    \item An emulated evaluation framework that models CXL load/store semantics on commodity dual-socket hardware.
    \item A head-to-head comparison of CHIME on RDMA vs.\ CXL for a subset of YCSB workloads.
    \item An analysis of which of CHIME's techniques (hopscotch leaves, vacancy-aware locks, metadata replication, sibling-based validation, speculative read) remain load-bearing when latency per op drops by~10$\times$.
\end{itemize}

%% ============================================================
\section{Background}
\label{sec:background}
%% ============================================================

\subsection{CHIME in One Page}
%% TODO: Brief recap from Part One — hybrid B+~tree / hopscotch leaves, six core
%% techniques, 10~CN + 1~MN disaggregated memory architecture.

\subsection{CXL Memory}
%% TODO: CXL.mem protocol, load/store semantics vs one-sided RDMA verbs,
%% ~200ns load latency vs ~2us RDMA READ. Why this matters for CHIME.

\subsection{Emulating CXL on Commodity Hardware}
%% TODO: CloudLab r650 has two NUMA nodes (one per socket). By pinning the
%% "client" threads to node 0 and binding memory to node 1, we emulate
%% cross-device load/store via QPI/UPI. Discuss the fidelity of this emulation
%% (same memory semantics, different absolute latency than real CXL).

%% ============================================================
\section{Part One Recap: RDMA Baseline}
\label{sec:part1-recap}
%% ============================================================

%% TODO: One-page summary of Part One results:
%%  - What reproduced (fig_12 C/D/E, fig_15a, fig_15b)
%%  - What did not (fig_12 A/B/LOAD — time/hardware constraints)
%%  - Where the baseline numbers live (cite Part One report)
For completeness, we briefly recap the RDMA reproduction results from the Part One progress report.
%% Figure reference: figures/fig_12_rdma.pdf  (copy from Part One)

%% ============================================================
\section{Port Design}
\label{sec:design}
%% ============================================================

\subsection{Architecture Overview}
%% TODO: Diagram — current RDMA path vs new CXL path.
%%  - User code → DSM → ThreadConnection → libibverbs → NIC → remote MN
%%  - User code → CxlDSM → CxlTransport → NUMA-bound memory → load/store

\subsection{The \texttt{CxlTransport} Abstraction}
%% TODO: Describe CxlTransport interface. Methods: read_sync, write_sync,
%% cas_sync, faa_boundary_sync, read_dm_sync, barrier, alloc, free.
%% All synchronous (no coroutine polling). Map GlobalAddress to a pointer
%% into the CXL-backed mmap region.

\subsection{\texttt{CxlDSM}: Drop-in DSM Replacement}
%% TODO: Explain CxlDSM inherits the public DSM API, so Tree.cpp is unchanged.
%% Show the preprocessor alias ``#define DSM CxlDSM'' trick and why it works.

\subsection{Build-System Surgery}
%% TODO: Explain the 5 blocker fixes committed in 152b4e9:
%%  1. Common.h gated Boost/RDMA includes behind #ifndef USE_CXL
%%  2. addr_from_ror corrected to use ga.val = ror.dest (was corrupting batch addresses)
%%  3. is_on_chip routing restored in batch operations
%%  4. DSM.cpp / RDMA-only sources filtered out of CXL build via CMake
%%  5. USE_CXL option() declaration moved before first reference
%% Include relevant diff snippets.

\subsection{Memory Allocation and Address Translation}
%% TODO: Single-node means GlobalAddress.nodeID is always 0. The .offset field
%% indexes into a contiguous mmap region. On-chip memory is emulated via a
%% small fixed-size region at the start of the allocation space.

\subsection{What We Did NOT Port}
%% TODO:
%%  - Multi-node support (single-node only; coordination via shared memory)
%%  - Coroutines (all operations synchronous; no CQ polling)
%%  - Memcached-based initialization (no cluster handshake needed)
%%  - RPC-based directory operations (rpc stubs return 0)

%% ============================================================
\section{Evaluation}
\label{sec:eval}
%% ============================================================

\subsection{Experimental Setup}
%% TODO: Hardware: r650 @ Clemson — dual-socket Intel Xeon Gold 6338N,
%% 256GB DDR4, two NUMA nodes. OS: Ubuntu 20.04. Build: GCC 9, CMake.
%% Workloads: YCSB C/D/E (reproduce from Part One).
%% Baselines:
%%  1. Single-node CHIME on RDMA (loopback, same node as both CN and MN)
%%  2. Single-node CHIME on CXL (NUMA node 0 threads, node 1 memory)
%%  3. Single-node Sherman, SMART for reference (if time permits)

\subsection{Throughput and Latency}
%% TODO: Figure placeholder — mirror fig_12 from paper.
%% Discuss how CXL's load/store semantics change the throughput-latency curve.

\subsection{Which CHIME Techniques Still Matter?}
%% TODO: With ~10x lower per-op latency, do hopscotch leaves still beat B+~tree
%% leaves? Does vacancy-aware locking still help when contention window shrinks?
%% Disable each technique in turn (mirror fig_15a) and measure the delta.

\subsection{Cache Consumption on CXL}
%% TODO: On CXL, the index cache optimization may be less important since
%% "remote" access is cheap. Measure and discuss.

%% ============================================================
\section{Discussion}
\label{sec:discussion}
%% ============================================================

\subsection{Does CHIME Still Make Sense on CXL?}
%% TODO: The paper's core claim is cache efficiency via hopscotch leaves.
%% On CXL where loads are fast, does this trade-off still favor hopscotch?
%% Speculate, then let data decide.

\subsection{Limitations of NUMA Emulation}
%% TODO:
%%  - Latency is QPI/UPI (~80ns), real CXL is closer to 200-300ns
%%  - Bandwidth symmetric; real CXL devices often bandwidth-limited
%%  - No real CXL.mem protocol overhead (credit flow control, etc.)

\subsection{Path to a Real CXL Deployment}
%% TODO: What would change if run on real CXL hardware (e.g., a CXL memory
%% expander attached via PCIe 5.0)? DMA APIs, persistence model, failure domain.

%% ============================================================
\section{Related Work}
\label{sec:related}
%% ============================================================

%% TODO: CXL-native index structures, disaggregated memory surveys,
%% other CXL port efforts (e.g., Pond, TPP).

%% ============================================================
\section{Conclusion}
\label{sec:conclusion}
%% ============================================================

%% TODO: Summary paragraph + 1 sentence on open artifact.

%% ============================================================
\bibliographystyle{plainnat}
\bibliography{references}
%% ============================================================

\appendix

\section{Reproducibility}
\label{app:repro}

%% TODO:
%%  - Repo: github.com/djjay0131/CHIME, branch: main, commit: [SHA]
%%  - Build: cmake -DUSE_CXL=ON .. && make -j
%%  - Run: exp/fig_12_cxl.py
%%  - Hardware: any dual-socket x86_64 with NUMA support
%%  - Results reproduced on: CloudLab r650 @ Clemson

\end{document}
